\section{The Architect's Question}\label{the-architects-question}

Before any model is selected, any GPU is counted, any latency budget is
allocated, the architect must answer one question: \emph{what is the
actual token-scale cost of delivering the stakeholder's request?} This
chapter exists to make that question automatic. We do not reach for
vendor marketing numbers or rule-of-thumb heuristics. Instead we trace a
vague stakeholder ask --- ``we need a Q\&A system over our internal
documents'' --- through a canonical loop that turns it into concrete
architectural bounds. The output is a set of derived quantities (token
throughput, KV-cache memory, prefill FLOPs, decode bandwidth) that are
verifiable, derived from first principles, and grounded in the canonical
\textasciitilde2,000 user, \textasciitilde10 rps / \textasciitilde40 rps
peak, RAG Q\&A, 70B dense FP16 scenario that every other chapter in this
handbook cites. We use a we-voice throughout: no shoulds, no musts. Only
``we'' --- the architect and the system --- and the reasoning that
connects them.

\section{Concept}\label{concept}

The unit we price, size, and optimise is the \textbf{token}. We
established what a token is in \hyperref[chap:1]{Chapter 1} --- not a word, character, or
byte, but whatever the model's tokenizer carves text into via learned
byte-pair merging; that token counts are never proportional to character
counts; and that ``one token ≈ 4 characters'' is a rough rule for
general-English prose, to be confirmed against the \emph{actual}
tokenizer before pricing anything {[}Ch. 1 §1{]}. We do not restate that
here; this chapter builds on it.

What \hyperref[chap:22]{Chapter 22} adds is not a re-derivation of the token but a
\textbf{re-framing of it as the input to a decision loop}. The architect
treats a token as a metered unit of thought --- the way a kilowatt-hour
meters electricity --- not because one token is meaningful by itself,
but because every downstream cost and capacity number is denominated in
it. The durable mental model, unchanged from \hyperref[chap:1]{Chapter 1}'s: \emph{tokens
are the interface between a human's request and a machine's arithmetic,
and anything we are asked to size or price reduces to ``how many tokens,
in what context, with what precision.''} The loop below turns that unit
into concrete bounds.

\section{Mental Model}\label{mental-model}

Think of inference as two gears turning at different speeds. The prefill
gear is large and toothed: it does a lot of work per revolution (many
FLOPs), but it only turns once per request. The decode gear is small and
smooth: it turns every token, doing relatively little work per step, but
it keeps turning for every output token. The prefill gear is limited by
how fast the compute engines can crank (FLOPS); the decode gear is
limited by how fast data can be fed from memory (GB/s). When we ask ``is
this workload compute-bound or bandwidth-bound?'' the answer depends
entirely on which gear is doing the work.

A useful concrete image: prefill is like filling a bucket from a fire
hose --- a burst of high-rate flow that fills the entire capacity in one
go. Decode is like drinking from that bucket one sip at a time --- each
sip is small, but we keep sipping until the bucket is empty. The hose is
limited by pipe width (bandwidth); the sips are limited by cup size and
swallowing rate (compute).

The architect's mental model, distilled: \textbf{tokens are the unit of
measure; prefill and decode are opposite-bound processes that require
different levers.} This distinction is the single biggest discriminator
between an architecture that meets its SLOs and one that does not.

\section{Worked Example --- The Canonical
Loop}\label{worked-example-the-canonical-loop}

We now demonstrate the canonical loop: turning a vague stakeholder ask
into concrete bounds. The stakeholder says: \emph{``We need a Q\&A
system over our internal documents. About 2,000 employees will use it,
they'll ask questions throughout the day, and the answers should be
accurate and fast.''} From this single sentence we will derive a full
set of token-scale quantities that every downstream chapter will price
against.

\textbf{Step 1 --- Identify the unit.} The ask mentions ``Q\&A over
internal documents.'' We do not assume ``one question = one page'' or
``one question = 500 words.'' We go to the model's tokenizer and ask:
how many tokens does a typical prompt contain? We extract representative
prompts from the documented corpus, run them through the tokenizer, and
measure. In the canonical scenario, the prompt consists of
\textasciitilde1,200 tokens of query text plus \textasciitilde8,000
tokens of retrieved RAG context, yielding \textbf{\textasciitilde9,200
input tokens}. The answer generated is \textasciitilde300 tokens. These
numbers are not assumed; they are measured {[}1P{]}.

\textbf{Step 2 --- Establish the traffic profile.} The ask says ``about
2,000 employees'' and ``throughout the day.'' We instrument or survey to
find the per-user rate. The canonical scenario settles on
\textasciitilde10 requests/s average, with peaks of \textasciitilde40
rps. These are derived quantities {[}2° DERIVED{]}, not stipulated.

\textbf{Step 3 --- Compute the token throughput.} With \(I = 9{,}200\)
input tokens and \(O = 300\) output tokens per request, and
\(\lambda = 10\) req/s average:

\[
\text{input tokens/s} = I \cdot \lambda = 9{,}200 \times 10 = 92{,}000 \text{ tokens/s}
\]

\[
\text{output tokens/s} = O \cdot \lambda = 300 \times 10 = 3{,}000 \text{ tokens/s}
\]

\[
\text{input/output ratio} = \frac{I}{O} = \frac{9{,}200}{300} \approx 30\times
\]

(At the \textasciitilde40 req/s peak these rise to
\textasciitilde368,000 and \textasciitilde12,000 tokens/s respectively.)
{[}2° DERIVED{]}

\textbf{Step 4 --- Derive architecture-relevant quantities.} The 30×
input/output ratio is the single most important derived quantity here.
It means this workload is input-heavy: the great majority of latency,
memory, and energy is spent in prefill (processing the 9.2K prompt), not
decode (generating the 300 answer tokens). This ratio alone dictates
that KV-cache memory and prefill compute dominate the design --- not
decode bandwidth. Every downstream chapter (memory, compute, cost) will
price against these derived numbers.

\textbf{Step 5 --- Record in \hyperref[tab:22.1]{Table 22-1}.} The full table of derived
quantities appears below.

This loop --- stakeholder ask → measured token counts → traffic profile
→ token throughput → architecture-relevant derived quantities --- is the
canonical thought process an AI solution architect runs, every time,
before a single architecture decision is made. It is how we move from
``we need a Q\&A system'' to ``our prefill will be compute-bound at
\textasciitilde1.29 PFLOP per request, and our decode will be
bandwidth-bound at \textasciitilde5.6 TB/s per token.'' The loop is the
chapter's central contribution; the quantities that issue from it are
the evidence {[}1P{]}/{[}2°{]}/{[}VERIFY{]} that every following chapter
trades on.

\protect\phantomsection\label{tab-22-1}
\subsection{\hyperref[tab:22.1]{Table 22-1} --- Derived quantities from the canonical loop
(worked example, not
reference)}\label{table-22-1-derived-quantities-from-the-canonical-loop-worked-example-not-reference}

\begin{longtable}[]{@{}
  >{\raggedright\arraybackslash}p{0.3333\linewidth}
  >{\raggedright\arraybackslash}p{0.3333\linewidth}
  >{\raggedright\arraybackslash}p{0.3333\linewidth}@{}}
\toprule\noalign{}
\begin{minipage}[b]{\linewidth}\raggedright
metric
\end{minipage} & \begin{minipage}[b]{\linewidth}\raggedright
value
\end{minipage} & \begin{minipage}[b]{\linewidth}\raggedright
derivation
\end{minipage} \\
\midrule\noalign{}
\endhead
\bottomrule\noalign{}
\endlastfoot
input tokens / request & \textasciitilde9,200 & 1,200 prompt + 8K RAG
context {[}1P DERIVED{]} \\
output tokens / request & \textasciitilde300 & generated answer length
{[}1P DERIVED{]} \\
input / output ratio & \textasciitilde30× & 9,200 ÷ 300 {[}2°
DERIVED{]} \\
input tokens/s @ 10 rps & \textasciitilde92,000 & 9,200 × 10 {[}2°
DERIVED{]} \\
input tokens/s @ peak 40 rps & \textasciitilde368,000 & 9,200 × 40 {[}2°
DERIVED{]} \\
output tokens/s @ 10 rps & \textasciitilde3,000 & 300 × 10 {[}2°
DERIVED{]} \\
output tokens/s @ peak 40 rps & \textasciitilde12,000 & 300 × 40 {[}2°
DERIVED{]} \\
prefill FLOPs per request & \textasciitilde1.29 PFLOP & 2 × 70B × 9.2K ≈
1.29 × 10¹⁵ {[}DERIVED{]} \\
decode bandwidth per token & \textasciitilde5.6 TB/s & 140 GB weights /
25 ms TPOT {[}DERIVED{]} \\
KV-cache memory per request & \textasciitilde12 GB (70B, 8-bit KV) & 1.3
MB/token × 9,200 tokens {[}1P DERIVED{]} \\

\label{tab:22.1}\end{longtable}

\begin{figure}
\centering
\pandocbounded{\includegraphics[keepaspectratio,alt={\hyperref[fig:22.1]{Fig 22.1} - Prefill vs decode bottleneck divergence {[}2° DERIVED{]}},width=\maxwidth{\textwidth},keepaspectratio]{figures/fig-22-2201.pdf}}
\caption{\hyperref[fig:22.1]{Fig 22.1} - Prefill vs decode bottleneck divergence {[}2°
DERIVED{]}}

\label{fig:22.1}\end{figure}

\emph{\hyperref[fig:22.1]{Fig 22.1} --- Prefill FLOPs grow sharply with context length
(linear 2NL, red) and faster still once the quadratic attention term is
added (orange band, +17\% at 9.2K → \textasciitilde2.4× at 128K); the
9.2K canonical point is marked. Prefill is compute-bound
(\textasciitilde1.29 PFLOP/request → \textasciitilde1.19 PFLOPS required
at the \textasciitilde1.08 s budget vs H100 \textasciitilde0.989 PFLOPS
peak). Decode is absent from this axis because it is a fixed
\textasciitilde140 GFLOP/token (\textasciitilde0.00014 PFLOP,
\textasciitilde9,000× below this axis' floor) and stays bandwidth-bound
(Ch6): the two regimes are deliberately drawn on different resources. }

\section{Measurement}\label{measurement}

For this chapter, measurement reduces to \textbf{counting tokens
correctly}, because a miscount at the input silently misprices
everything downstream. Three things an architect can and should check:

\begin{enumerate}
\def\labelenumi{\arabic{enumi}.}
\item
  \textbf{Actual token count, not the rule of thumb.} Run the model's
  \emph{own tokenizer} on representative prompts from the actual
  traffic. The ``4 chars ≈ 1 token'' heuristic is for estimation only;
  real counts differ by language, formatting, code, and tokenizer
  version. Measure on our traffic, not the marketing figure.
\item
  \textbf{Input vs output split.} Measure both legs of the request
  (prompt tokens and generated tokens), because they land on different
  bottlenecks --- input on memory/prefill, output on decode --- and on
  different cost line items. A request that is 90\% input and 10\%
  output has a very different cost profile than one that is 50/50, even
  at the same total token count.
\item
  \textbf{Peak vs average context.} Log the distribution of context
  lengths, not just the mean. A workload that averages 9.2K input tokens
  but peaks at 32K (illustrative variant) has a very different KV-cache
  and latency profile. The 90th-percentile context length is often the
  number the architect must design for.
\end{enumerate}

This measurement habit is the token-layer answer to the book's recurring
question, ``what would I actually measure here?'' We measure token
counts and their distribution, at the edge, before any architecture
decision is made.

\section{Common Mistakes}\label{common-mistakes}

\begin{itemize}
\item
  \textbf{Assuming one token ≈ one word.} It is a useful heuristic for
  general English prose and nothing more. Code, numerics, and non-Latin
  scripts routinely run 2--5× the naive estimate, which misprices
  capacity and cost.
\item
  \textbf{Ignoring the input/output asymmetry.} Treating a request as
  ``one unit'' hides that input-heavy workloads are memory-bound and
  output-heavy ones are decode-bound --- opposite bottlenecks with
  opposite fixes. The 30× input/output ratio in the canonical loop is
  the concrete illustration of this principle.
\item
  \textbf{Quoting context window as free capacity.} The window is an
  upper bound, not a recommendation; using it fully is expensive.
  Keeping a 9.2K average inside a 128K window still costs for the length
  we actually use.
\item
  \textbf{Trusting a vendor's tokenizer count without checking.}
  Tokenizer versions change and report differently; measure on
  \emph{our} traffic, not the marketing figure.
\item
  \textbf{Skipping the canonical loop.} Assuming that the stakeholder's
  description is sufficient to size the system. The loop exists because
  vague asks almost always underspecify the token-scale cost; skipping
  it produces architectures that fail latency or cost SLOs.
\end{itemize}

\section{Architecture Consequence}\label{architecture-consequence}

The canonical loop's derived quantities have immediate and concrete
architecture consequences. The 30× input/output ratio means this
workload's prefill phase dominates: \textasciitilde1.29 PFLOP of compute
per request, and a KV-cache that grows with 9.2K input tokens. The
decode phase, while smaller in total tokens, requires \textasciitilde5.6
TB/s of HBM bandwidth per generated token --- a bandwidth-bound design.

These opposite bottlenecks lead to a central architectural decision:
\textbf{can we disaggregate prefill and decode?} If prefill needs FLOPS
and decode needs bandwidth, a single homogeneous GPU pool is suboptimal.
A two-pool architecture --- a prefill cluster optimized for compute
(more GPUs, higher FLOPS, model parallelism) and a decode cluster
optimized for bandwidth (faster HBM, PagedAttention, continuous
batching) --- can achieve the same serving SLO with fewer total GPUs
than a homogeneous design. This is the central theme of \hyperref[chap:11]{Chapter 11}
(Serving) and Pattern 4 (Prefill/decode disaggregation). The architect
who runs the canonical loop and records the derived quantities in \hyperref[tab:22.1]{Table 22-1} is already positioned to make this decision with numbers, not
intuition.

Importantly, the loop does not end at a serving topology; it ends at a
\emph{decision about where the intelligence lives}. A 70B dense pool on
the fleet is one placement; an 8-bit or MoE variant, a retrieval-first
design that buys capability with context rather than parameters, a
test-time-search loop that spends compute on hard queries, or an agent
runtime that orchestrates several specialised models are all alternative
placements of the same capability. The canonical loop's derived
quantities --- token throughput, KV footprint, prefill FLOPs, TCO ---
are precisely the numbers an architect uses to compare those placements.
So the full hierarchy the loop supports is:

\begin{quote}
Requirement → Workload → Intelligence (where the capability comes from)
→ System → Architecture → Fleet → \textbf{Decision}.
\end{quote}

The \textbf{Decision} rung is where the loop's output is committed ---
an ADR, an owner, a review trigger (Chapters 25--26). Every earlier
derived quantity exists to inform that final, committed decision, and
Appendix A's closing question --- \emph{where should the intelligence
live?} --- is the architectural form of that same question, made
explicit at the top of the hierarchy rather than hidden at the bottom.

\section{What We Still Don't Know}\label{what-we-still-dont-know}

This chapter's numbers lean on the same open questions the book raises
at the mechanism layer. Rather than restating them, we point to the
fuller treatment:

\begin{itemize}
\tightlist
\item
  \textbf{Sustained HBM bandwidth for weight-read kernels} --- discussed
  in \hyperref[chap:2]{Chapter 2} §7; the 3.35 TB/s H100 figure is a theoretical peak and
  real serving kernels may sustain less. {[}VERIFY HYPOTHESIS{]}
\item
  \textbf{Prefill FLOP cost under continuous batching} --- \hyperref[chap:2]{Chapter 2} §7;
  the 2 × params × tokens approximation ignores activation reuse across
  batched requests. {[}VERIFY HYPOTHESIS{]}
\item
  \textbf{Effect of quantization on decode vs prefill} --- \hyperref[chap:2]{Chapter 2} §7;
  quantization shifts both bottlenecks but the precise trade-off point
  is workload-dependent. {[}VERIFY HYPOTHESIS{]}
\item
  \textbf{Cross-technology bandwidth numbers} --- \hyperref[chap:2]{Chapter 2} §7; HBM3e
  and MI300X publish higher peaks, which move the GPU-count equation
  without changing the compute-bound vs bandwidth-bound classification.
  {[}2° FACT{]}
\end{itemize}

\section{End-of-Chapter Mini-Case}\label{end-of-chapter-mini-case}

The architect is still in the early design conversation about the
internal Q\&A tool described throughout this handbook. The team has
settled on a 70B-class dense model for accuracy, and they have a rough
traffic estimate: \textasciitilde2,000 registered employees, each
roughly as active as the canonical scenario (\textasciitilde10
requests/s average, peaks \textasciitilde40 rps), with prompts of
\textasciitilde1,200 tokens of internal documents plus \textasciitilde8K
retrieved context and \textasciitilde300 tokens of answer.

From the token layer alone, the architect can already state the hard
numbers: each request carries \textasciitilde9.2K input tokens and
\textasciitilde300 output tokens. At 10 rps average, the system must
sustain \textasciitilde92,000 input tokens/s in prefill and
\textasciitilde3,000 output tokens/s in decode. At 40 rps peak, those
numbers jump to \textasciitilde368,000 and \textasciitilde12,000
respectively. The architect can also state the two opposite bottlenecks:
prefill will be compute-bound (\textasciitilde1.29 PFLOP per request at
9.2K input), and decode will be bandwidth-bound (\textasciitilde5.6 TB/s
required weight-read rate per token). The team's immediate architectural
choice is whether to build a single homogeneous GPU cluster or to
separate prefill and decode pools --- the arithmetic from this chapter
makes clear that a homogeneous design will be suboptimal for either
bottleneck, and that the disaggregation option explored in \hyperref[chap:11]{Chapter 11} is
worth evaluating early.

Before passing this workload to \hyperref[chap:23]{Chapter 23} (Memory), the architect locks
in one more number: the tokens-per-request split. This is the currency
every downstream chapter will price against, and it has already been
established as \textasciitilde9.2K input + \textasciitilde300 output per
request. The rest of the design --- model fitting, memory budget,
latency budget, cost --- can now proceed in token-denominated terms.
